% Generated draft: requires mathematical and source review.
\documentclass[11pt]{article}
\usepackage[margin=25mm]{geometry}
\usepackage{fontspec}
\setmainfont{Latin Modern Roman}
\usepackage{amsmath,amssymb,mathtools,unicode-math}
\setmathfont{Latin Modern Math}
\usepackage{enumitem,needspace,xurl,hyperref}
\setlength{\parindent}{0pt}
\setlength{\parskip}{5pt}
\setlength{\emergencystretch}{3em}
\newcommand{\N}{\mathbb{N}}
\newcommand{\Z}{\mathbb{Z}}
\newcommand{\Q}{\mathbb{Q}}
\newcommand{\R}{\mathbb{R}}
\newcommand{\C}{\mathbb{C}}
\newcommand{\F}{\mathbb{F}}
\newcommand{\A}{\mathbb{A}}
\newcommand{\PP}{\mathbb{P}}
\newcommand{\E}{\mathbb{E}}
\newcommand{\eps}{\varepsilon}
\newcommand{\dd}{\,\mathrm{d}}
\begin{document}
\section*{UnsolvedMath ID: 11000001}
\textbf{Source problem code:} catalogue entry
\subsection{Definitions and mathematical statement}
For a genus \(g\) in the source's range, let \(\mathcal M_g\) denote its moduli space of genus-\(g\) curves/Riemann surfaces and let
\[\Psi:\mathcal M_g\longrightarrow[-\log(\mu_g),\infty)\]
be the particular function introduced in the omitted preceding text. The parameter \(\mu_g\) must be positive for the stated real logarithm. The codomain states a lower bound and does not by itself assert that the lower bound is attained or that \(\Psi\) is surjective.

The programme has three tasks: study this \(\Psi\) as a \emph{rational Morse function}; classify its rational critical points; and derive properties of \(H^*(\mathcal M_g;\Q)\), including the specific statements referred to later in the source. Neither the formula for \(\Psi\), \(\mu_g\), the meaning of rational Morse/critical point nor the later statements are supplied. In particular ``rational critical point'' cannot silently be interpreted as an algebraic point defined over \(\Q\), or as an ordinary smooth critical point, without that definition.

If the intended framework is ordinary orbifold Morse theory, one would analyze local lifts of \(\Psi\), their critical loci, Hessian indices and isotropy, and explain the resulting filtration/attachments computing rational cohomology. This describes one conditional framework, not an assertion that the given function is smooth, proper or Morse, nor a replacement of the source's rational notion. The moduli space has finite-automorphism/orbifold issues and is noncompact in the usual closed-surface setting, so these features must enter any claimed cohomological deduction. No universal cohomology presentation or stable-range result is added as a substitute for this function-specific programme.
\subsection{Short English statement}
Study the specified moduli-space function through its intended rational Morse theory, classify its critical points and extract rational cohomological consequences. The defining function and rational terminology are missing from this excerpt.
\subsection{Sources}
\begin{itemize}
\item Local UnsolvedMath registry, permanent ID \(11000001\), exact source code \texttt{catalogue entry}.
\item Farb, \emph{Problems on Mapping Class Groups and Related Topics}, Problem 1.1, PDF page 11: \url{https://www.math.uchicago.edu/~farb/papers/mcgbook.pdf}. Registry access 2026-07-29; pdf-font-structure extraction. The summary attributes this programme to Farb (2006).
\item The summary names Harer and Madsen--Weiss for other rational-cohomology/stable-range progress, without exact citations; it does not claim they resolve the specified Morse programme.
\end{itemize}
\end{document}
